\appendix
\setcounter{figure}{0}
\renewcommand{\thefigure}{A\arabic{figure}}
\renewcommand{\theHfigure}{A.\arabic{figure}}
\renewcommand{\theHequation}{A.\arabic{equation}}
\section{Viewing Geometry and Single-Night Trajectories}
\label{app:checks}

\subsection{Solar elongation and angular rate}
\label{app:geometry}

Solar elongation is calculated from target and solar directions
with the same observing origin. In J2000-equatorial coordinates,
let $\boldsymbol r_{\rm o}(t)$ be the observer's heliocentric
position and $\boldsymbol r_{\rm a}(t-\tau)$ the asteroid's
light-time-corrected heliocentric position. The unit directions
and elongation are
\begin{align}
\hat{\boldsymbol u}_{\rm a} &=
\frac{\boldsymbol r_{\rm a}(t-\tau)-\boldsymbol r_{\rm o}(t)}
 {|\boldsymbol r_{\rm a}(t-\tau)-\boldsymbol r_{\rm o}(t)|},\\
\hat{\boldsymbol u}_{\odot} &=
-\frac{\boldsymbol r_{\rm o}(t)}{|\boldsymbol r_{\rm o}(t)|},\\
\epsilon &=\cos^{-1}\left(\hat{\boldsymbol u}_{\rm a}\cdot
 \hat{\boldsymbol u}_{\odot}\right).
\end{align}
The elongation distribution in Figure~\ref{fig:geometry} counts
each primary association once.

Angular rates use adjacent associations of the same identity
within one observing night, in different exposures and separated
by a positive time interval. The rate is the great-circle
separation divided by the elapsed time. There are \yrNRatePairs\
such pairs; their median predicted and measured rates are
\yrPredictedRate\ and \yrMeasuredRate~\rateunit, respectively.
The two rates are calculated independently from the predicted
and measured coordinates.

\begin{figure}[htbp]
\centering
\annualplot{figA1_motion_and_solar_elongation.pdf}{\linewidth}
\caption{First-year viewing geometry and apparent motion.
(a) Predicted and measured angular rates from adjacent exposures
of the same known identity within one night.
(b) Solar elongation of the primary known-asteroid associations,
using a common observing origin for the Sun and asteroid.
The rate distributions count exposure pairs; the elongation
distribution counts individual associations.}
\label{fig:geometry}
\end{figure}

\subsection{Single-night trajectory fits}
\label{app:shortarc}

The short-arc calculation constructs Lambert endpoints from
equal heliocentric-distance hypotheses with zero assumed radial
derivatives. The provisional tracklet-subset search admits seeds
with final-epoch heliocentric speeds below 140~km~s$^{-1}$.
The subsequent refit uses the 35~km~s$^{-1}$ acceptance limit
specified in Section~\ref{sec:motion_method}. Positions and
trajectory spans are evaluated at the exposure midpoints.

Figure~\ref{fig:shortarc} compares the observed spans with the
angular fit RMS of the \yrNCandidates\ retained candidates.
The C/2025~Y1 links are identified separately. The RMS measures
agreement along the observed short arc. Longer baselines
provide the additional constraints on distance and radial
motion needed for orbit determination.

\begin{figure}[htbp]
\centering
\annualplot{figA2_candidate_short_arc_diagnostic.pdf}{0.75\linewidth}
\caption{Exposure-midpoint span and short-arc fit RMS for the
retained first-year single-night candidates. Links associated
with C/2025~Y1 are shown separately. The RMS describes the
angular agreement of the trajectory over the measured arc.}
\label{fig:shortarc}
\end{figure}
